% preamble.tex — định dạng chung của khóa luận.
% Quy ước (theo khóa luận khóa trước cùng GVHD): Times New Roman, hình đánh số "Hình 3.1",
% chú thích in nghiêng ở DƯỚI hình, sau mỗi hình có đoạn diễn giải; bảng đánh số "Bảng 3.1", chú thích ở TRÊN.

% ---------------- Phông chữ và cỡ chữ ----------------
\usepackage{fontspec}
% TeX Gyre Termes: bản tương đương Times New Roman, đủ dấu tiếng Việt, có sẵn trên Overleaf.
\setmainfont{TeX Gyre Termes}
\setmonofont{DejaVu Sans Mono}[Scale=0.85]
% Chữ thân bài 13pt
\makeatletter
\renewcommand\normalsize{\@setfontsize\normalsize{13pt}{16pt}}
\makeatother
\normalsize

% ---------------- Trang và giãn dòng ----------------
\usepackage[a4paper,top=2.5cm,bottom=2.5cm,left=3cm,right=2cm]{geometry}
\usepackage{setspace}
\setstretch{1.3}
\setlength{\parindent}{1cm}
\setlength{\parskip}{4pt}
\usepackage{indentfirst}

% ---------------- Bảng, hình, danh sách ----------------
\usepackage{graphicx}
\graphicspath{{images/}}
\usepackage{float}
\usepackage{booktabs}
\usepackage{tabularx}
\usepackage{longtable}
\usepackage{array}
\usepackage{multirow}
\usepackage[table]{xcolor}
\usepackage{enumitem}
\setlist{nosep,leftmargin=1.2cm}
\renewcommand{\arraystretch}{1.25}
\newcolumntype{L}{>{\raggedright\arraybackslash}X}

% Hình: chú thích nghiêng, ở dưới; bảng: chú thích ở trên. Đánh số theo chương.
\usepackage{caption}
\captionsetup{labelsep=period,justification=centering,font=normalsize}
\captionsetup[figure]{name=Hình,font=it,position=below,skip=6pt}
\captionsetup[table]{name=Bảng,position=above,skip=4pt}

% ---------------- Tiêu đề chương, mục ----------------
\usepackage{titlesec}
\titleformat{\chapter}[display]
  {\centering\bfseries\fontsize{16}{20}\selectfont}
  {CHƯƠNG \thechapter}{6pt}{\MakeUppercase}
\titlespacing*{\chapter}{0pt}{0pt}{18pt}
\titleformat{\section}{\bfseries\fontsize{14}{18}\selectfont}{\thesection.}{0.5em}{}
\titleformat{\subsection}{\bfseries\fontsize{13}{16}\selectfont}{\thesubsection.}{0.5em}{}
\titleformat{\subsubsection}{\bfseries\itshape\fontsize{13}{16}\selectfont}{\thesubsubsection.}{0.5em}{}
\setcounter{secnumdepth}{3}
\setcounter{tocdepth}{2}

% Tên các danh mục bằng tiếng Việt
\renewcommand{\contentsname}{MỤC LỤC}
\renewcommand{\listfigurename}{DANH MỤC HÌNH ẢNH}
\renewcommand{\listtablename}{DANH MỤC BẢNG}
\renewcommand{\bibname}{TÀI LIỆU THAM KHẢO}
\renewcommand{\chaptername}{Chương}

% Chương không đánh số (Mở đầu, Kết luận...) vẫn có trong mục lục
\newcommand{\chuongkhongso}[1]{%
  \chapter*{\MakeUppercase{#1}}%
  \addcontentsline{toc}{chapter}{\MakeUppercase{#1}}}

% ---------------- Liên kết, đầu trang ----------------
\usepackage[hidelinks,unicode]{hyperref}
\usepackage{fancyhdr}
\pagestyle{fancy}
\fancyhf{}
\fancyfoot[C]{\thepage}
\renewcommand{\headrulewidth}{0pt}
\fancypagestyle{plain}{\fancyhf{}\fancyfoot[C]{\thepage}\renewcommand{\headrulewidth}{0pt}}

% ---------------- Tiện ích viết bài ----------------
% Đánh dấu chỗ chưa viết (hiện chữ xám nghiêng trong PDF, tìm lại bằng Ctrl+F "chuaviet")
\newcommand{\chuaviet}[1][Nội dung sẽ được bổ sung ở giai đoạn sau.]{%
  \par\noindent{\color{gray}\itshape (#1)}\par}
% Tên file, đường dẫn, mã trong câu
\newcommand{\code}[1]{\texttt{#1}}
\setlength{\emergencystretch}{3em}
